\chapter{Kommentierte Quellen}

Die folgenden achtundzwanzig Werke sind nicht bloß genannt, sondern eingeordnet: Zu jedem steht, wofür man es aufschlägt und was es leistet, das andere nicht leisten.

\lit{1}{Aigner, M., Ziegler, G.: \textit{Das BUCH der Beweise}. 6.\ Aufl., Springer 2018.}{Eine Sammlung besonders schöner Beweise, mehrere davon graphentheoretisch --- darunter Turán in fünf Fassungen und der Satz von Cayley in vier. Wer sehen will, dass es zu einem Ergebnis verschiedene Wege gibt, findet hier das beste Anschauungsmaterial. Deutsch verfügbar.}

\lit{2}{Alon, N., Spencer, J.: \textit{The Probabilistic Method}. 4.\ Aufl., Wiley 2016.}{\textbf{Das Standardwerk zur Methode aus Kapitel 13 und 14.} Beginnt mit Erdős' Ramsey-Schranke --- genau der Rechnung aus Aufgabe 13.3 --- und führt bis zu den heutigen Verfeinerungen. Der Anhang über Konzentrationsungleichungen ist eigenständig brauchbar.}

\lit{3}{Appel, K., Haken, W.: \textit{Every Planar Map is Four Colorable}. AMS 1989.}{Die vollständige Darstellung des Vierfarbenbeweises von 1976, einschließlich der Konfigurationslisten. Historisch bedeutsam, zum Lernen ungeeignet --- der Wert liegt darin, den Umfang der Fallprüfung mit eigenen Augen zu sehen.}

\lit{4}{Berge, C.: \textit{Théorie des graphes et ses applications}. Dunod 1958.}{Das Buch, das die Graphentheorie in Frankreich begründete und den Satz aus Kapitel 11 über vergrößernde Wege enthält. Bergescher Stil: knapp, elegant, gelegentlich zu knapp. Von historischem Interesse und noch immer lesbar.}

\lit{5}{Bollobás, B.: \textit{Modern Graph Theory}. Springer 1998.}{Deckt mehr Gebiete ab als jedes andere Lehrbuch dieses Umfangs --- Zufallsgraphen, algebraische Methoden, extremale Fragen, Tutte-Polynom. Anspruchsvoll geschrieben; als zweites Buch nach einem Einstieg sehr gut, als erstes zu steil.}

\lit{6}{Bondy, J.\,A., Murty, U.\,S.\,R.: \textit{Graph Theory}. Springer 2008.}{Umfangreich, mit vielen Anwendungen und einem ausführlichen Aufgabenteil. Die Beweise sind vollständiger ausgeschrieben als bei Diestel, dafür ist die Auswahl weniger streng. Gute Wahl, wenn man beim Lesen ungern nachdenkt, was ausgelassen wurde.}

\lit{7}{Diestel, R.: \textit{Graphentheorie}. 5.\ Aufl., Springer 2017.}{\textbf{Das beste Lehrbuch des Gebiets, und die Quelle für alles, was dieser Band nicht bewiesen hat.} Kuratowski in Kapitel 4, Brooks in Kapitel 5, Tutte in Kapitel 2. Deutsch und englisch verfügbar; die englische Fassung ist frei zugänglich. Knapp, aber lückenlos --- Diestel schreibt nie \enquote{man sieht leicht}, wo etwas zu zeigen ist.}

\lit{8}{Erdős, P., Rényi, A.: \textit{On the Evolution of Random Graphs}. Publ.\ Math.\ Inst.\ Hung.\ Acad.\ Sci.\ 5, 1960.}{Die Gründungsarbeit der Zufallsgraphentheorie und der Ort, an dem der Phasenübergang aus Kapitel 14 zuerst beschrieben wird. Vierzig Seiten, gut lesbar, frei verfügbar --- eine der wenigen klassischen Arbeiten, die man ohne Vorbereitung durchgehen kann.}

\lit{9}{Euler, L.: \textit{Solutio problematis ad geometriam situs pertinentis}. Comm.\ Acad.\ Sci.\ Petropol.\ 8, 1741 (vorgetragen 1736).}{Die Geburtsurkunde des Gebiets. Bemerkenswert ist, wie wenig Euler von dem benutzt, was wir heute für nötig halten: kein Graph, keine Kante, kein Grad --- nur eine Buchführung über Brücken und ein Paritätsargument. Lateinisch; Übersetzungen sind leicht zu finden.}

\lit{10}{Ford, L.\,R., Fulkerson, D.\,R.: \textit{Flows in Networks}. Princeton 1962.}{Die Buchfassung des Satzes aus Kapitel 12, mit vielen Folgerungen --- darunter Menger, Hall und Kőnig als Sonderfälle, genau wie in diesem Band. Die Darstellung ist bis heute nicht überholt.}

\lit{11}{Frieze, A., Karoński, M.: \textit{Introduction to Random Graphs}. Cambridge 2016.}{Modernes Lehrbuch zu Kapitel 14, mit vollständigen Beweisen der Schwellenwertsätze --- auch der zweiten Momentenmethode, die dieser Band nur zitiert hat. Frei als Vorabfassung verfügbar.}

\lit{12}{Godsil, C., Royle, G.: \textit{Algebraic Graph Theory}. Springer 2001.}{\textbf{Für Kapitel 14 die richtige Fortsetzung.} Behandelt Spektren, Automorphismengruppen und stark reguläre Graphen. Der Petersen-Graph tritt hier in fast jedem Kapitel auf --- man versteht nach der Lektüre, warum er ein Sonderfall ist.}

\lit{13}{Gonthier, G.: \textit{Formal Proof --- The Four-Color Theorem}. Notices of the AMS 55(11), 2008.}{Elf Seiten darüber, was es heißt, einen Beweis vollständig zu formalisieren. Der Aufsatz ist auch für Leser ohne Vorkenntnisse in Beweisassistenten verständlich und die beste kurze Einführung in die Debatte aus Kapitel 9.}

\lit{14}{Graham, R., Rothschild, B., Spencer, J.: \textit{Ramsey Theory}. 2.\ Aufl., Wiley 1990.}{Das Standardwerk zu dem Gebiet, das Kapitel 13 nur angetippt hat. Enthält die Sätze von van der Waerden, Hales--Jewett und Szemerédi. Anspruchsvoll, aber mit sorgfältigen Beweisen.}

\lit{15}{Hall, P.: \textit{On Representatives of Subsets}. J.\ London Math.\ Soc.\ 10, 1935.}{Drei Seiten. Der Satz aus Kapitel 11, in der ursprünglichen Fassung über Repräsentantensysteme --- ohne Heiratseinkleidung, ohne Graphen. Lesenswert als Beispiel dafür, wie viel sich in wenigen Seiten sagen lässt.}

\lit{16}{Harary, F.: \textit{Graph Theory}. Addison-Wesley 1969.}{Das Lehrbuch, das die englische Terminologie festgelegt hat. Heute inhaltlich überholt, aber als Nachschlagewerk für Begriffsgeschichte nützlich --- wer wissen will, warum eine Sache so heißt, findet hier meist die Antwort.}

\lit{17}{Kőnig, D.: \textit{Theorie der endlichen und un\-end\-li\-chen Graphen}. Akad.\ Verlagsges.\ 1936.}{\textbf{Das erste Lehrbuch der Graphentheorie überhaupt}, auf Deutsch geschrieben. Enthält die Sätze aus Kapitel 2 und 11 in ihrer ursprünglichen Fassung. Die deutsche Grundbegrifflichkeit dieses Bandes --- Ecke, Kante, Grad --- geht darauf zurück.}

\lit{18}{Kuratowski, K.: \textit{Sur le problème des courbes gauches en topologie}. Fund.\ Math.\ 15, 1930.}{Die Arbeit hinter Satz 8.8. Kuratowski behandelt das Problem topologisch, nicht kombinatorisch --- was erklärt, warum der Beweis so viel schwerer ist, als die Aussage vermuten lässt.}

\lit{19}{Lovász, L.: \textit{Combinatorial Problems and Exercises}. 2.\ Aufl., AMS 2007.}{Über sechshundert Aufgaben mit Lösungen, nach Schwierigkeit gestaffelt. Wer diesen Band durchgearbeitet hat und weiterüben will, findet hier den natürlichen Anschluss. Die Lösungen sind knapp, aber vollständig.}

\lit{20}{Menger, K.: \textit{Zur allgemeinen Kurventheorie}. Fund.\ Math.\ 10, 1927.}{Der Satz aus Kapitel 12 --- in einer Arbeit über Kurventheorie, ohne jeden Bezug zu Flüssen oder Netzwerken. Ein Musterbeispiel dafür, dass mathematische Ergebnisse ihre Anwendungen oft um Jahrzehnte vorwegnehmen.}

\lit{21}{Mohar, B., Thomassen, C.: \textit{Graphs on Surfaces}. Johns Hopkins 2001.}{Die Fortsetzung von Kapitel 8 über die Ebene hinaus: Einbettungen in Torus, Kleinsche Flasche und allgemeine Flächen. Enthält auch, was über die verbotenen Minoren des Torus bekannt ist --- und was nicht.}

\lit{22}{Nešetřil, J., Ossona de Mendez, P.: \textit{Sparsity}. Springer 2012.}{Eine Theorie dünner Gra\-phen\-klas\-sen, die das Gebiet in den letzten fünfzehn Jahren geprägt hat. Anspruchsvoll; der Bezug zu diesem Band liegt bei den Kapiteln 8 und 13.}

\lit{23}{Robertson, N., Seymour, P.: \textit{Graph Minors I--XX}. J.\ Combin.\ Theory Ser.\ B, 1983--2004.}{\textbf{Der längste Beweis der Graphentheorie}, über zwanzig Arbeiten und mehr als fünfhundert Seiten. Niemand liest ihn am Stück. Die Arbeiten V und XIII enthalten die zentralen Strukturaussagen; ein Überblick von Lovász (2005) ist der bessere Einstieg.}

\lit{24}{Schrijver, A.: \textit{Combinatorial Optimization}. Springer 2003.}{Drei Bände, über sechzehnhundert Seiten. Die maßgebliche Darstellung zu Paarungen, Flüssen und Matroiden --- alles aus den Kapiteln 11 und 12 in voller Allgemeinheit, mit umfangreichen historischen Anmerkungen, die den Geschichtskästen dieses Bandes zugrunde liegen.}

\lit{25}{Spielman, D.: \textit{Spectral and Algebraic Graph Theory}. Vorlesungsskript, Yale, 2019.}{\textbf{Frei verfügbar und für Kapitel 14 die beste Fortsetzung.} Behandelt Laplace-Spektren, Expander und das spektrale Zerlegen von Netzwerken. Spielmans Stil ist ungewöhnlich zugänglich für ein Gebiet dieser Art.}

\lit{26}{Turán, P.: \textit{Egy gráfelméleti szélsőértékfeladatról}. Mat.\ Fiz.\ Lapok 48, 1941.}{Die Arbeit aus dem Zwangsarbeitslager, ungarisch. Eine deutsche Fassung erschien 1954 unter dem Titel \kursiv{Über ein Extremalproblem der Graphentheorie}. Sieben Seiten, die ein ganzes Teilgebiet begründeten.}

\lit{27}{West, D.: \textit{Introduction to Graph Theory}. 2.\ Aufl., Prentice Hall 2001.}{Der amerikanische Standard für Einführungsvorlesungen. Ausführlicher als Diestel, mit sehr vielen Aufgaben und einer sorgfältigen Behandlung der Beweistechniken. Wenn dieser Band zu knapp war, ist West der nächste Schritt.}

\lit{28}{Wilson, R.: \textit{Four Colors Suffice}. Erweiterte Aufl., Princeton 2013.}{Die Geschichte des Vier\-far\-ben\-pro\-blems über $124$ Jahre, für ein all\-ge\-mei\-nes Publikum geschrieben und dennoch mathematisch genau. Enthält Kempes fehlerhaften Beweis samt Heawoods Gegenbeispiel --- die Vorlage für Aufgabe 9.3.}

\begin{merksatz}
\boxtitel{Wenn Sie nur drei Bücher kaufen}
\begin{center}
\renewcommand{\arraystretch}{1.6}
\begin{tabular}{@{}P{1.0cm}P{2.8cm}P{4.4cm}@{}}
\toprule
\textbf{[7]} & \textbf{Diestel} & \textbf{Für die Theorie}\\
\textbf{[27]} & \textbf{West} & Für die Übung\\
\textbf{[2]} & \textbf{Alon, Spencer} & \textbf{Für die Methode}\\
\bottomrule
\end{tabular}
\end{center}

\begin{center}
\begin{minipage}{0.88\textwidth}
\kursiv{Diestel ist das Buch, zu dem man zurückkehrt --- knapp, streng, vollständig. West ist das Buch, mit dem man arbeitet --- ausführlich, aufgabenreich, geduldig. Alon und Spencer sind das Buch, das die Denkweise ändert: Nach dem ersten Kapitel sieht man Existenzfragen anders als vorher.}
\end{minipage}
\end{center}

\begin{center}
\begin{minipage}{0.88\textwidth}
\kursiv{Zwei Ergänzungen, beide frei verfügbar: [25] Spielman, wenn Kapitel 14 Sie interessiert hat, und [8] Erdős/Rényi --- vierzig Seiten, die man an einem Abend liest und danach anders über Zufall denkt.}
\end{minipage}
\end{center}
\end{merksatz}

\bigskip
\begin{flushright}
\itshape\color{grau}
Vier Landstücke, sieben Brücken,\\
und ein Spaziergang, den es nicht gibt.\\[10pt]
\normalfont\color{grau}\upshape
--- der Anfang dieses Bandes\\
\itshape und sein Maßstab:\\
Nichts ist selbstverständlich.\\
Jeder Schritt wird erklärt.
\end{flushright}
